\documentclass[11pt]{article}
\usepackage[utf8]{inputenc}
\usepackage[T1]{fontenc}
\usepackage{newtxtext,newtxmath}
\usepackage{amsmath}
\usepackage{multicol}
\usepackage{geometry}
\usepackage{enumitem}
\usepackage{xcolor}
\usepackage{titlesec}

% Configurações de layout
\geometry{a4paper, left=1cm, right=1cm, top=0.5cm, bottom=1.2cm}
\setlength{\columnseprule}{0.4pt}
\setlength{\baselineskip}{1.0\baselineskip}

% Cores para títulos
\titleformat{\section}{\normalfont\Large\color{blue}}{\thesection}{1em}{}
\titleformat{\subsection}{\normalfont\large\color{red}}{\thesubsection}{1em}{}

\title{\textcolor{blue}{Questionário: \textit{Men Against Fire} (Black Mirror)}}
\author{Professor: Jefferson}
\date{}

\begin{document}

\maketitle
\vspace{-1cm}

\begin{multicols}{2}

\section*{Orientações}
Leia atentamente o texto de apoio e responda \textbf{no seu caderno}. Use exemplos do episódio sempre que possível para justificar suas respostas.

\section*{Texto de Apoio}

O episódio \textit{Men Against Fire} (Engenharia Reversa), o quinto da terceira temporada de \textit{Black Mirror}, apresenta um futuro onde soldados são equipados com um implante neural chamado "MÁSCARA" (MASS). Este sistema otimiza o desempenho militar, fornecendo dados táticos em tempo real, mas também altera drasticamente a percepção sensorial dos combatentes. O protagonista, Stripe, participa de missões para exterminar as "baratas", seres descritos como mutantes perigosos que ameaçam a pureza genética da humanidade.

A trama se aprofunda quando o implante de Stripe começa a falhar após um confronto. Ele passa a enxergar a realidade sem os filtros tecnológicos, descobrindo que as "baratas" são, na verdade, seres humanos comuns, perseguidos por uma ideologia de eugenia e limpeza social. O episódio explora como a tecnologia pode ser usada para desumanizar o "outro", facilitando atos de violência extrema ao remover a empatia, o remorso e a culpa dos agressores.

\textit{Men Against Fire} serve como uma metáfora poderosa para a propaganda de guerra e a manipulação da percepção pública. Ao transformar o inimigo em um monstro literal através da realidade aumentada, o sistema garante que os soldados cumpram suas ordens sem questionamentos morais. O cenário levanta discussões éticas sobre o uso da inteligência artificial no campo de batalha e as consequências psicológicas da desumanização institucionalizada.

\section*{Parte 1: Compreensão do Episódio}

\hspace{0.6cm}1. Qual é a função principal do implante "MÁSCARA" (MASS) nos soldados?

2. Como os soldados visualizam as "baratas" enquanto o implante está funcionando perfeitamente?

3. O que acontece com Stripe durante a missão no vilarejo que causa o mau funcionamento de seu implante?

4. Qual é a verdadeira identidade das "baratas" revelada após a falha do sistema de Stripe?

5. Além da visão, como o sistema MÁSCARA altera os sentidos do olfato e da audição dos soldados?

6. Qual é a justificativa "científica" ou social dada pelo governo para o extermínio das "baratas"?

7. O que Arquette (o psicólogo militar) revela sobre a dificuldade histórica de fazer soldados matarem em guerras reais?

8. Como a população civil local (os moradores das vilas) trata as "baratas" e por que eles colaboram com os militares?

9. Qual é o dilema moral oferecido a Stripe quando ele é confrontado com a verdade no centro de detenção?

10. Descreva a cena final do episódio e explique o que a visão de Stripe representa em contraste com a realidade.

\section*{Parte 2: Análise Crítica}

\hspace{0.6cm}11. Explique como a desumanização do inimigo é utilizada como ferramenta de controle no episódio.

12. O episódio aborda o conceito de eugenia. De que forma o sistema justifica o extermínio de uma parte da população em nome de um "bem maior"?

13. Como a manipulação dos sentidos (visão, audição e olfato) remove a barreira da empatia humana?

14. Faça um paralelo entre a tecnologia MÁSCARA e o uso de propaganda de ódio ou *fake news* em conflitos contemporâneos.

15. Você acredita que a tecnologia pode ser usada para silenciar a consciência moral de um indivíduo? Justifique sua resposta.

\section*{Parte 3: Conexão Pessoal}

\hspace{0.6cm}16. Se você vivesse nessa sociedade e descobrisse a verdade sobre o sistema, o que faria?

17. De que forma as "bolhas" das redes sociais hoje podem funcionar como uma "MÁSCARA" moderna, distorcendo nossa visão sobre quem pensa diferente de nós?

18. Como podemos exercitar a empatia em um mundo que frequentemente nos incentiva a ver o "outro" como um inimigo ou alguém inferior?

19. Qual a importância de questionar as informações e a realidade que nos são apresentadas por autoridades ou tecnologias?

\end{multicols}

\end{document}
